% !PRINTED 602
\weberpage{602}
%<<CONT
haben wir die Discriminante der Zahl $\omega$ genannt. Das Verh\"altniss $D:d$ ist eine ganze rationale Quadratzahl. Die beiden Wurzeln von (1) erhalten den Ausdruck

\weq{3}{\omega = \frac{b+\sqrt{D}}{2c}, \quad \omega' = \frac{b-\sqrt{D}}{2c}.}
Setzen wir
\weq{4}{D = e^2 d,}
so wird
\weq{5}{\omega = \frac{b+e\sqrt{d}}{2c}, \quad \omega' = \frac{b-e\sqrt{d}}{2c},}
und wegen (1) sind dies dann und nur dann ganze Zahlen, wenn $c=1$ ist.

Nun ist $D\equiv 0$ oder $\equiv 1\ (\mathrm{mod}\ 4)$, und wenn also $D$ gerade ist, so ist es durch 4 theilbar; $b$ ist in diesem Falle gerade, und weil $d$ nicht durch 4 theilbar sein kann, so muss auch $e$ gerade sein.

Ist aber $D$ ungerade, so sind auch $b$ und $e$ ungerade. Dieser Fall kann aber, wie (4) zeigt, nur dann eintreten, wenn $d\equiv 1\ (\mathrm{mod}\ 4)$ ist.

Daraus erhalten wir folgendes Resultat:

1. \gesperrt{Ist $d\equiv 2$ oder $\equiv 3\ (\mathrm{mod}\ 4)$, so ist die Zahl}
\[ \omega = x + y\sqrt{d} \]
\gesperrt{dann und nur dann eine ganze Zahl, wenn die rationalen Zahlen $x, y$ ganze Zahlen sind.} (Quadratische Irrationalzahlen erster Art, nach \gesperrt{Dirichlet}.)

2. \gesperrt{Ist $d\equiv 1\ (\mathrm{mod}\ 4)$, so ist}
\[ \omega = \frac{x+y\sqrt{d}}{2} \]
\gesperrt{dann und nur dann eine ganze Zahl, wenn die rationalen Zahlen $x,y$ entweder zwei gerade oder zwei ungerade ganze Zahlen sind.} (Quadratische Irrationalzahlen erster oder zweiter Art, je nachdem $e$ gerade oder ungerade ist.)

Letzteres k\"onnen wir auch so ausdr\"ucken:

3. \gesperrt{Ist $d\equiv 1\ (\mathrm{mod}\ 4)$, so ist}
\[ \omega = x + y\,\frac{1+\sqrt{d}}{2} \]
\gesperrt{dann und nur dann eine ganze Zahl, wenn die rationalen Zahlen $x$ und $y$ ganze Zahlen sind.}
